\clearpage
\chapter*{Errata e complemento à exposição n\textsuperscript{o} 236}
\addcontentsline{toc}{chapter}{Errata e complemento à exposição
n\textsuperscript{o} 236}
\src{303}

\noindent\textbf{8. TDTE VI: Os esquemas de Picard. Propriedades gerais.}

\begin{center}
[Seminário Bourbaki, t. 14, 1961/62, n\textsuperscript{o} 236, 23 p.]
\end{center}

\noindent\textbf{Página 12, linha 16.} --- Em vez de «reta afim», leia-se
«reta afim privada da origem».

\noindent\textbf{Linha 18.} --- Em vez de «$X[t]$», leia-se
«$X[t,t^{-1}]$».

\noindent\textbf{Página 16, proposição 3.1, 6\textsuperscript{a} linha.} ---
Em vez de «Ora, a hipótese ...», leia-se «Ora, da hipótese ...».

\noindent\textbf{Página 21, § 4: O teorema de finitude para o esquema de
Picard.} --- As questões de finitude do gênero das levantadas neste parágrafo
foram resolvidas quase totalmente desde a redação da presente exposição.
Indiquemos os principais fatos agora conhecidos nessa direção. Para simplificar
os enunciados, supomos implicitamente que todos os pré-esquemas de Picard que
neles intervêm existem, embora uma modificação evidente desses enunciados
permita eliminar toda hipótese de existência. Doravante, $S$ designa um esquema
noetheriano, e $X,Y$, esquemas próprios sobre $S$.

\medskip
\noindent(i) Seja $f:X\longrightarrow Y$ um morfismo sobrejetivo; então
\[
f^*: \operatorname{Pic}_{Y/S}\longrightarrow\operatorname{Pic}_{X/S}
\]
é um morfismo de tipo finito.
\src{304}
\noindent(ii) Suponhamos $Y$ projetivo sobre $S$, munido de um Módulo
invertível amplo relativamente a $S$, e seja $X$ o esquema dos zeros de uma
seção qualquer desse Módulo, $f:X\longrightarrow Y$ a imersão canônica.
Suponhamos, enfim, que as componentes irredutíveis das fibras de $Y/S$ tenham
dimensão $\geq3$. Então
\[
f^*: \operatorname{Pic}_{Y/S}\longrightarrow\operatorname{Pic}_{X/S}
\]
é de tipo finito.

\noindent(iii) Suponhamos que $X$ seja projetivo sobre $S$ e que todas as suas
fibras geométricas sejam íntegras e de dimensão $n$. Seja
$\mathcal O_X(1)$ um Módulo invertível amplo sobre $X$, que permita definir
polinômios de Hilbert. Seja $M$ um subconjunto de
$\operatorname{Pic}_{X/S}$; então $M$ é quase-compacto se, e somente se, nos
polinômios de Hilbert
\[
a_0x^n+a_1x^{n-1}+a_2x^{n-2}+\cdots
\]
dos elementos de $M$, os coeficientes $a_1$ e $a_2$ permanecem limitados, omitindo-se todo coeficiente cujo monômio teria grau negativo; assim, apenas $a_1$ intervém se $n=1$, e nenhum dos dois se $n=0$.

\noindent(iv) Para todo inteiro $n\neq0$, o homomorfismo de elevação à
$n$-ésima potência no pré-esquema em grupos $\operatorname{Pic}_{X/S}$ é um
morfismo de tipo finito.

Notemos que (i) e (ii) significam também que (sob as hipóteses consideradas
nesses enunciados (i) e (ii)) um subconjunto $M$ de
$\operatorname{Pic}_{Y/S}$ é quase-compacto se, e somente se, sua imagem em
$\operatorname{Pic}_{X/S}$ o é. Conclui-se daí que um Módulo invertível
$\mathcal L$ sobre $Y$ é $\tau$-equivalente a zero se, e somente se, sua imagem
inversa sobre $X$ o é; em outros termos, $\operatorname{Pic}_{Y/S}^\tau$ é a
imagem inversa de $\operatorname{Pic}_{X/S}^\tau$. Em particular, para mostrar
que o primeiro pré-esquema é de tipo finito sobre $S$, basta prová-lo para o
segundo, pois $\operatorname{Pic}_{Y/S}\longrightarrow
\operatorname{Pic}_{X/S}$ é de tipo finito. Utilizando então (i), o lema de
Chow e (iii), obtemos:

\medskip
\noindent(v) $\operatorname{Pic}_{X/S}^\tau$ é de tipo finito sobre $S$.

De modo geral, a conjunção de (i) para um morfismo finito e de (ii) permite
reduzir, para a maioria das questões de finitude, ao caso em que $X/S$ tem
fibras geométricas íntegras normais de dimensão $\leq2$; muitas vezes, aplicando
(i) a um morfismo sobrejetivo não necessariamente finito e a resolução das
singularidades das superfícies algébricas (devida, em característica qualquer,
a Abhyankar), reduzimos ao caso em que $X/S$ é até mesmo liso, portanto tem
fibras geométricas não singulares de dimensão 2. Isso permite, por exemplo,
tendo em conta (v) e a desigualdade de Picard--Igusa que majora o posto do grupo
de Néron--Severi de uma superfície projetiva não singular, provar a seguinte
generalização do teorema de finitude de Néron:

\medskip
\noindent(vi) Seja $X/S$ próprio sobre $S$ e, de resto, arbitrário; então os
grupos de Néron--Severi
\[
\operatorname{Pic}_{X_i/k_i}/\operatorname{Pic}_{X_i/k_i}^\circ
\]
das fibras geométricas $X_i/k_i$ de $X/S$
\src{305}
são de tipo finito, e seu posto e a ordem de seu subgrupo de torção permanecem
limitados.

O mesmo método de redução ao caso das superfícies não singulares e os teoremas
conhecidos nesse caso (a saber, o teorema de finitude de Néron e o fato de que a
forma de interseção sobre o grupo de Néron--Severi é não degenerada) implicam:

\medskip
\noindent(vii) Seja $X$ um esquema próprio sobre um corpo algebricamente
fechado. Existe um número finito de curvas fechadas íntegras $C_i$
($1\leq i\leq r$) em $X$, com a seguinte propriedade: para todo subconjunto
$M$ de $\operatorname{Pic}_{X/k}$, $M$ é quase-compacto se, e somente se, os
inteiros
\[
\deg_{C_i'}\mathcal L\qquad(\mathcal L\in M)
\]
permanecem limitados (onde $C_i'$ designa a normalização de $C_i$).

Pode-se tomar aqui, para $r$, o posto do grupo de Néron--Severi. Uma vez sabido
que este último é de tipo finito, (vii) reduz-se ao fato de que as formas
lineares sobre o grupo de Néron--Severi definidas por curvas $C$ em $X$ só se
anulam simultaneamente sobre os elementos de torção do grupo de Néron--Severi.
No caso em que $X$ é projetivo e não singular, esse resultado, assim como (v),
era devido a Matsusaka. Utilizando (vii), obtém-se facilmente a seguinte
caracterização dos Módulos invertíveis $\tau$-equivalentes a zero sobre $X$:

\medskip
\noindent(viii) Sejam $X/k$ um esquema próprio sobre um corpo e $\mathcal L$ um
Módulo invertível sobre $X$. Condições equivalentes:

\noindent a. $\mathcal L$ é $\tau$-equivalente a zero;

\noindent b. Para todo Módulo coerente $F$, tem-se
\[
\chi(F\otimes\mathcal L)=\chi(F),
\]
onde $\chi$ designa a característica de Euler--Poincaré;

\noindent b'. Como (b), mas com $F=\mathcal O_Y$, sendo $Y$ um subesquema
fechado íntegro de dimensão 1 em $X$;

\noindent c. Para todo $Y$ como acima, designando por $Y'$ a curva normalizada,
tem-se
\[
\deg\mathcal L_{Y'}=0.
\]

Se $X/k$ é projetivo e munido de um Módulo invertível amplo
$\mathcal O_X(1)$, as condições precedentes são também equivalentes às
seguintes:

\noindent d. Para todo inteiro $m$, $\mathcal L^{\otimes m}(1)$ é amplo.

\noindent e. (Se $X$ é íntegro.) Para todo par de inteiros $m,n$, tem-se
\[
\chi(\mathcal L^{\otimes m}(n))=\chi(\mathcal O_X(n))
\]
\src{306}
(isto é, (b) é verdadeiro tomando $F=\mathcal L^{\otimes m}(n)$).

Para a suficiência desta última condição, note-se que ela significa que os
polinômios de Hilbert dos $\mathcal L^{\otimes m}$ são todos iguais; portanto,
em virtude do critério de Mumford (iii), os $\mathcal L^{\otimes m}$ permanecem
em um subconjunto quase-compacto de $\operatorname{Pic}_{X/k}$, isto é, vale
(a). As condições (b), (b') e (c) devem ser consideradas variantes (sobre um
esquema próprio qualquer) da noção de equivalência numérica, definida
habitualmente sobre variedades projetivas não singulares. Para estas últimas,
a equivalência de (a) e (c) era evidentemente bem conhecida (Matsusaka).

O critério (e) implica também o resultado seguinte:

\medskip
\noindent(ix) Seja $f:X\longrightarrow S$ um morfismo projetivo e plano cujas
fibras geométricas são íntegras. Então $\operatorname{Pic}_{X/S}^\tau$ é aberto
e fechado em $\operatorname{Pic}_{X/S}$.

Limitar-nos-emos a alguns comentários sobre a demonstração dos resultados-chave
(i), (ii), (iii) (o resultado (iv) é um pouco apartado dos demais e se demonstra
utilizando somente (i) para os morfismos finitos sobrejetivos radiciais, mais
precisamente, para um morfismo de Frobenius). Para (i), utilizam-se de modo
essencial as ideias da descida não plana (ver, em particular, TDTE I, página 9).
Ocorre que (tendo em vista apenas resultados de finitude) a falta de critérios
de efetividade para dados de descida é inofensiva. Por outro lado, Mumford
demonstrou recentemente uma forma um pouco menos forte de (iii), a saber, o
critério que envolve todos os coeficientes do polinômio de Hilbert. Seu
argumento é extremamente simples e inspirado pela demonstração de um critério
de amplitude de Nakai (enunciado por este último para as superfícies não
singulares e generalizado por Mumford aos morfismos projetivos quaisquer).
Parece-me, aliás, que esse argumento só é válido mediante uma ligeira restrição
adicional sobre as fibras de $X/S$, a propriedade $(S_2)$ de Serre, verificada
se as fibras geométricas são normais. Utiliza-se então esse critério restrito na
demonstração de (ii): o critério (i) permite, com efeito, reduzir ao caso em que
$Y/S$ é plano com fibras geométricas íntegras e normais e, aplicando o critério
de Mumford, reduz-se facilmente ao caso em que $X/S$ satisfaz as mesmas
condições. Da hipótese de dimensão resulta então que as fibras geométricas de
$Y$ e de $X$ têm profundidade $\geq2$ em seus pontos fechados, o que permite
aplicar os «critérios de equivalência» sob a forma que lhes é dada em SGA 1962
(exposição XII) e concluir a demonstração de (ii). Uma vez que dispomos de (i) e
(ii), não é difícil eliminar do critério de Mumford toda hipótese de normalidade
sobre as fibras e demonstrá-lo também sob a forma mais forte dada em (iii).
\src{307}
Notemos, enfim, que a demonstração de (i) e (ii) mostra igualmente que, no
caso em que $S$ é o espectro de um corpo $k$, o morfismo
\[
\operatorname{Pic}_{Y/k}\longrightarrow\operatorname{Pic}_{X/k}
\]
é afim (e não apenas de tipo finito).
